\section{Tạo cây thư mục giống project}

Project F450 được tổ chức theo kiểu một ROS 2 workspace, bên trong có một
package mô tả robot tên là \texttt{f450\_description}. Phần mô hình robot,
controller, script chạy trong Isaac Sim và script vẽ đồ thị được đặt chung
trong package này.

\subsection{Cấu trúc thư mục chính}

Cấu trúc rút gọn của project:

\begin{lstlisting}[language=bash]
f450/
|-- f450.usd
|-- f450_ver2.usd
|-- f450_ver3.usd
|-- f450_ver3_1.usd
|-- configuration/
|   |-- f450_ver3_base.usd
|   |-- f450_ver3_physics.usd
|   |-- f450_ver3_robot.usd
|   `-- f450_ver3_sensor.usd
|-- docs/
|   |-- Baocao.tex
|   |-- Baocao.pdf
|   |-- figures/
|   `-- sections/
`-- ros2_ws/
    |-- build/
    |-- install/
    |-- log/
    `-- src/
        `-- f450_description/
            |-- CMakeLists.txt
            |-- package.xml
            |-- urdf/
            |   `-- f450.urdf.xacro
            |-- meshes/
            |   |-- bottom_plate.stl
            |   |-- up_plate.stl
            |   |-- leg.stl
            |   |-- leg_mount.stl
            |   `-- protectors.stl
            `-- src/
                |-- f450_controller/
                |-- scripr_editor/
                |-- data/
                `-- test/
\end{lstlisting}

Trong đó:

\begin{itemize}
    \item \texttt{urdf/}: chứa file mô tả robot.
    \item \texttt{meshes/}: chứa các file STL được URDF tham chiếu.
    \item \texttt{src/f450\_controller/}: chứa toàn bộ controller Python.
    \item \texttt{src/scripr\_editor/}: chứa các đoạn script dùng trong Isaac
    Sim Script Editor.
    \item \texttt{src/data/}: chứa CSV log và các script vẽ đồ thị.
    \item \texttt{configuration/}: chứa các USD cấu hình được tách theo robot,
    physics và sensor.
\end{itemize}

\subsection{Tạo workspace và package}

Nếu tạo lại từ đầu, có thể bắt đầu bằng cây thư mục như sau:

\begin{lstlisting}[language=bash]
mkdir -p f450/ros2_ws/src
cd f450/ros2_ws/src

ros2 pkg create f450_description --build-type ament_cmake

cd f450_description
mkdir -p urdf meshes launch
mkdir -p src/f450_controller
mkdir -p src/scripr_editor
mkdir -p src/data
mkdir -p src/test
\end{lstlisting}

Sau đó đặt các file vào đúng vị trí:

\begin{lstlisting}[language=bash]
# Robot description
cp f450.urdf.xacro urdf/

# Mesh files used by the URDF
cp *.stl meshes/

# Python controller modules
cp *.py src/f450_controller/

# Isaac Sim Script Editor snippets
cp run.py set_pos.py circle_trajectory.py circle_run.py src/scripr_editor/

# Plotting and logging scripts
cp plot_tracking.py live_plot_tracking.py plot_trajectory_3d.py src/data/
\end{lstlisting}

Trong \texttt{CMakeLists.txt}, project cần install các thư mục mô tả robot:

\begin{lstlisting}[language=bash]
install(
  DIRECTORY urdf meshes launch
  DESTINATION share/${PROJECT_NAME}
)
\end{lstlisting}

Sau khi chuẩn bị xong, build workspace:

\begin{lstlisting}[language=bash]
cd f450/ros2_ws
colcon build
source install/setup.bash
\end{lstlisting}

\section{Import file URDF vào Isaac Sim}

File mô tả robot nằm trong thư mục:

\begin{lstlisting}[language=bash]
f450/ros2_ws/src/f450_description/urdf/f450.urdf.xacro
\end{lstlisting}

Isaac Sim URDF Importer thường làm việc trực tiếp với file \texttt{.urdf}. Nếu
file hiện tại là \texttt{.urdf.xacro}, cần generate ra \texttt{.urdf} trước:

\begin{lstlisting}[language=bash]
cd f450/ros2_ws
source install/setup.bash

ros2 run xacro xacro \
  src/f450_description/urdf/f450.urdf.xacro \
  -o src/f450_description/urdf/f450.urdf
\end{lstlisting}

Sau đó trong Isaac Sim:

\begin{enumerate}
    \item Mở Isaac Sim.
    \item Bật extension URDF Importer nếu chưa bật. Có thể tìm
    \texttt{omni.importer.urdf} trong cửa sổ Extensions.
    \item Mở công cụ URDF Importer.
    \item Chọn file \texttt{f450.urdf}.
    \item Kiểm tra đường dẫn mesh để các file STL trong thư mục
    \texttt{meshes/} được load đúng.
    \item Import robot vào stage.
    \item Lưu stage thành file USD, ví dụ \texttt{f450.usd} hoặc
    \texttt{f450\_ver3.usd}.
\end{enumerate}

Sau khi import, cần kiểm tra path của rigid body chính. Trong controller hiện
tại, path mặc định là:

\begin{lstlisting}[language=Python]
base_link_path = "/f450_simple/base_link"
\end{lstlisting}

Nếu trong stage Isaac Sim tên prim khác, cần sửa \texttt{base\_link\_path} khi
tạo controller.

\section{Viết controller ngoài và gọi từ Script Editor}

Controller được viết bên ngoài Isaac Sim trong thư mục:

\begin{lstlisting}[language=bash]
f450/ros2_ws/src/f450_description/src/f450_controller
\end{lstlisting}

Thư mục này là một Python package vì có file:

\begin{lstlisting}[language=bash]
f450_controller/__init__.py
\end{lstlisting}

Isaac Sim Script Editor không tự biết đường dẫn project, nên khi chạy cần thêm
đường dẫn controller vào \texttt{sys.path}. File \texttt{run.py} trong project
đang làm đúng việc này.

\subsection{Setup path và import controller}

Đoạn đầu tiên cần chạy trong Script Editor:

\begin{lstlisting}[language=Python]
import sys
import importlib
import os

CONTROLLER_PATH = "/config/Desktop/IssacSim_TA/f450/ros2_ws/src/f450_description/src"
SCRIPT_PATH = CONTROLLER_PATH + "/scripr_editor"
DATA_DIR = CONTROLLER_PATH + "/data"
TRACKING_CSV = DATA_DIR + "/f450_tracking.csv"

for path in (CONTROLLER_PATH, SCRIPT_PATH):
    if path not in sys.path:
        sys.path.insert(0, path)

os.makedirs(DATA_DIR, exist_ok=True)
importlib.invalidate_caches()
\end{lstlisting}

Nếu project nằm ở đường dẫn khác, chỉ cần sửa biến \texttt{CONTROLLER\_PATH}.

\subsection{Reload module khi sửa code}

Khi đang mở Isaac Sim, nếu sửa file controller bên ngoài, nên reload lại module
trước khi tạo controller mới:

\begin{lstlisting}[language=Python]
try:
    f450_app.stop()
except Exception:
    pass

MODULE_NAMES = [
    "f450_controller.control_utils",
    "f450_controller.motor_model",
    "f450_controller.altitude_hold",
    "f450_controller.attitude_pid",
    "f450_controller.position_hold",
    "f450_controller.yaw_hold",
    "f450_controller.motor_mixer",
    "f450_controller.propeller_spinner",
    "f450_controller.tracking_logger",
    "f450_controller.attitude_hold_compat",
    "f450_controller.issac_attitude_hold",
]

modules = {}
for module_name in MODULE_NAMES:
    modules[module_name] = importlib.import_module(module_name)

for module_name in MODULE_NAMES:
    modules[module_name] = importlib.reload(modules[module_name])

issac_attitude_hold = modules["f450_controller.issac_attitude_hold"]
\end{lstlisting}

Việc reload giúp Script Editor dùng phiên bản code mới nhất mà không cần tắt
mở lại Isaac Sim.

\subsection{Tạo controller và bắt đầu chạy}

Sau khi import, tạo đối tượng controller:

\begin{lstlisting}[language=Python]
f450_app = issac_attitude_hold.F450AttitudeHold(
    base_link_path="/f450_simple/base_link",
    z_target=1.0,
    pwm_hover=1650.0,
)

f450_app.position_angle_limit_deg = 6.5
f450_app.position_accel_limit = 1.2

f450_app.set_yaw_target(0.0)
f450_app.start_tracking_log(TRACKING_CSV, sample_period=0.02)
f450_app.start()

print("F450 controller STARTED")
print("Tracking CSV:", TRACKING_CSV)
\end{lstlisting}

Khi gọi \texttt{f450\_app.start()}, controller đăng ký callback
\texttt{on\_physics\_step(dt)} với physics engine của Isaac Sim. Từ thời điểm
đó, mỗi physics step sẽ đọc trạng thái drone, tính PID, mix PWM, cập nhật motor
model và apply lực vào body.

\section{Mở, chạy mô phỏng, log data và vẽ đồ thị}

\subsection{Mở mô phỏng}

Quy trình chạy mô phỏng:

\begin{enumerate}
    \item Mở Isaac Sim.
    \item Mở stage USD đã import robot, ví dụ \texttt{f450.usd} hoặc
    \texttt{f450\_ver3.usd}.
    \item Mở Script Editor.
    \item Paste hoặc mở nội dung file \texttt{src/scripr\_editor/run.py}.
    \item Kiểm tra lại \texttt{CONTROLLER\_PATH} và
    \texttt{base\_link\_path}.
    \item Nhấn Run trong Script Editor.
\end{enumerate}

Nếu mọi thứ đúng, terminal/console sẽ in:

\begin{lstlisting}[language=bash]
F450 controller STARTED
Tracking CSV: /config/Desktop/IssacSim_TA/f450/ros2_ws/src/f450_description/src/data/f450_tracking.csv
\end{lstlisting}

\subsection{Đặt setpoint tĩnh}

Để đổi target vị trí ngang, chạy trong Script Editor:

\begin{lstlisting}[language=Python]
f450_app.set_position_target(5.0, -1.0)
\end{lstlisting}

Lệnh này đặt:

\begin{equation}
    x_d=5.0,\quad y_d=-1.0
\end{equation}

Để đổi đồng thời $x,y,z$:

\begin{lstlisting}[language=Python]
f450_app.set_xyz_target(
    x_target=2.0,
    y_target=1.0,
    z_target=1.5,
)
\end{lstlisting}

Để chỉ đổi độ cao:

\begin{lstlisting}[language=Python]
f450_app.set_target(1.5)
\end{lstlisting}

Để đổi yaw target:

\begin{lstlisting}[language=Python]
f450_app.set_yaw_target(0.0)
\end{lstlisting}

\subsection{Bay theo quỹ đạo tròn}

Project có file:

\begin{lstlisting}[language=bash]
src/scripr_editor/circle_trajectory.py
src/scripr_editor/circle_run.py
\end{lstlisting}

Sau khi đã chạy \texttt{run.py} để tạo \texttt{f450\_app}, có thể chạy:

\begin{lstlisting}[language=Python]
import circle_trajectory as ct
import importlib
importlib.reload(ct)

ct.CIRCLE_RADIUS = 2.0
ct.CIRCLE_Z = 1.5
ct.CIRCLE_PERIOD = 20.0
ct.CIRCLE_CENTER_X = 0.0
ct.CIRCLE_CENTER_Y = 0.0
ct.TAKEOFF_TIME = 5.0
ct.LOOKAHEAD_DT = 0.18

ct.TRACKING_CSV = TRACKING_CSV
ct.start(f450_app)
\end{lstlisting}

Khi muốn dừng quỹ đạo tròn và lưu log:

\begin{lstlisting}[language=Python]
import circle_trajectory as ct
ct.stop(f450_app)
\end{lstlisting}

\subsection{Bật và dừng log}

Để bật log tracking:

\begin{lstlisting}[language=Python]
DATA_DIR = CONTROLLER_PATH + "/data"
TRACKING_CSV = DATA_DIR + "/f450_tracking.csv"
f450_app.start_tracking_log(TRACKING_CSV, sample_period=0.02)
\end{lstlisting}

Với \texttt{sample\_period=0.02}, logger ghi dữ liệu tối đa khoảng:

\begin{equation}
    f_s = \frac{1}{0.02} = 50\ \mathrm{Hz}
\end{equation}

Để dừng log:

\begin{lstlisting}[language=Python]
f450_app.stop_tracking_log()
\end{lstlisting}

File CSV sẽ chứa các nhóm cột chính:

\begin{itemize}
    \item Position: \texttt{x}, \texttt{x\_target}, \texttt{y},
    \texttt{y\_target}, \texttt{z}, \texttt{z\_target}.
    \item Attitude degree: \texttt{roll\_deg},
    \texttt{roll\_target\_deg}, \texttt{pitch\_deg},
    \texttt{pitch\_target\_deg}, \texttt{yaw\_deg},
    \texttt{yaw\_target\_deg}.
    \item Attitude radian: \texttt{roll\_rad},
    \texttt{roll\_target\_rad}, \texttt{pitch\_rad},
    \texttt{pitch\_target\_rad}, \texttt{yaw\_rad},
    \texttt{yaw\_target\_rad}.
\end{itemize}

\subsection{Vẽ đồ thị realtime}

Trong lúc Isaac Sim đang chạy, mở terminal khác và chạy:

\begin{lstlisting}[language=bash]
python3 ~/Desktop/IssacSim_TA/f450/ros2_ws/src/f450_description/src/data/live_plot_tracking.py \
  ~/Desktop/IssacSim_TA/f450/ros2_ws/src/f450_description/src/data/f450_tracking.csv \
  --wait \
  --max-points 2000
\end{lstlisting}

Option \texttt{--wait} cho phép script đợi tới khi file CSV xuất hiện. Option
\texttt{--max-points} giới hạn số điểm hiển thị để đồ thị realtime không quá
nặng.

\subsection{Vẽ đồ thị 2D sau mô phỏng}

Sau khi dừng mô phỏng và gọi \texttt{stop\_tracking\_log()}, tạo đồ thị 2D:

\begin{lstlisting}[language=bash]
python3 ~/Desktop/IssacSim_TA/f450/ros2_ws/src/f450_description/src/data/plot_tracking.py \
  ~/Desktop/IssacSim_TA/f450/ros2_ws/src/f450_description/src/data/f450_tracking.csv
\end{lstlisting}

Đồ thị 2D dùng để so sánh target và giá trị thực của:

\begin{equation}
    x,\ y,\ z,\ roll,\ pitch,\ yaw
\end{equation}

Vị trí dùng đơn vị mét, còn góc được vẽ theo độ.

\subsection{Vẽ quỹ đạo 3D}

Để xem quỹ đạo bay trong không gian 3D sau mô phỏng:

\begin{lstlisting}[language=bash]
python3 ~/Desktop/IssacSim_TA/f450/ros2_ws/src/f450_description/src/data/plot_trajectory_3d.py \
  ~/Desktop/IssacSim_TA/f450/ros2_ws/src/f450_description/src/data/f450_tracking.csv \
  --show
\end{lstlisting}

Để xem live 3D trong khi Isaac Sim đang chạy:

\begin{lstlisting}[language=bash]
python3 ~/Desktop/IssacSim_TA/f450/ros2_ws/src/f450_description/src/data/plot_trajectory_3d.py \
  ~/Desktop/IssacSim_TA/f450/ros2_ws/src/f450_description/src/data/f450_tracking.csv \
  --live --wait --interval 500
\end{lstlisting}

Để vẽ thêm vector attitude dọc theo đường bay:

\begin{lstlisting}[language=bash]
python3 plot_trajectory_3d.py f450_tracking.csv --show --attitude --attitude-step 15
\end{lstlisting}

Mặc định output ảnh sẽ có dạng:

\begin{lstlisting}[language=bash]
f450_tracking_tracking.png
f450_tracking_traj3d.png
\end{lstlisting}

\subsection{Dừng controller}

Khi kết thúc thử nghiệm, nên dừng log và controller:

\begin{lstlisting}[language=Python]
f450_app.stop_tracking_log()
f450_app.stop()
\end{lstlisting}

Nếu muốn chạy lại từ đầu, chạy lại file \texttt{run.py}. Trong file này đã có
đoạn:

\begin{lstlisting}[language=Python]
try:
    f450_app.stop()
except Exception:
    pass
\end{lstlisting}

Nhờ đó controller cũ được dừng trước khi tạo controller mới, tránh trường hợp
nhiều callback physics chạy chồng lên nhau.
